% !TEX root = ../main.tex
% ======================================================================
% Section 3 – Revue de littérature
% Responsable : Gaston Raphaël Kapawa
% Grille : pas de critère propre ; appuie « Cadrage & problématique » (6 pts)
% et « Rédaction & normes » (6 pts : au moins 15 sources, dont au moins 5
% articles scientifiques évalués par les pairs).
% Longueur indicative : ~2 pages.
% ======================================================================
% TODO : faire l'état des connaissances qui fonde l'analyse critique des sections 5 à 8.
% TODO : chaque source citée va dans shared/references.bib ET dans docs/sources-a-verifier.md.
% TODO : privilégier les articles évalués par les pairs (au moins 5 dans le rapport).

\subsection{Apprentissage fédéré en santé}
\label{sec:revue-fl-sante}
% TODO : origines de l'apprentissage fédéré (\ac{FedAvg}) et applications multicentriques en santé.
\arediger{apprentissage fédéré en santé}

\subsection{Attaques sur les mises à jour et confidentialité différentielle}
\label{sec:revue-attaques-dp}
% TODO : inférence d'appartenance, inversion de modèle, fuites par les gradients ; apport et limites de la \ac{DP}.
\arediger{attaques et confidentialité différentielle}

\subsection{Gouvernance des données et souveraineté autochtone}
\label{sec:revue-gouvernance}
% TODO : principes \pcap{}, \eptc{} (chapitre 9), \loivingtcinq{} ; travaux sur la gouvernance des modèles (et non seulement des données).
\arediger{gouvernance et souveraineté des données}

\subsection{Synthèse et lacunes}
\label{sec:revue-synthese}
% TODO : ce que la littérature ne règle pas, et qui justifie l'analyse des sections suivantes.
\arediger{synthèse et lacunes}
